\chapter{Planteamiento del problema}

La interpretación de resonancias magnéticas lumbares requiere revisar secuencias, identificar estructuras anatómicas, comparar niveles y, cuando corresponde, realizar mediciones. Aunque los visores médicos actuales permiten observar imágenes y efectuar mediciones manuales, la delimitación de estructuras y la organización de resultados cuantitativos pueden seguir dependiendo de tareas manuales o semimanuales. En RM lumbar, tanto en el plano sagital como en el axial, estructuras como cuerpos vertebrales, discos intervertebrales y canal espinal son relevantes para el análisis estructural. Su delimitación manual puede consumir tiempo, variar entre observadores y generar resultados difíciles de reutilizar si no quedan asociados a una estructura, un nivel anatómico y una imagen de origen. Esta situación afecta especialmente la trazabilidad de mediciones y observaciones cuando se trabaja con capturas, registros dispersos o informes redactados sin una salida estructurada.

El problema central que aborda este proyecto puede formularse de la siguiente manera:

Existe una oportunidad de asistencia informática en el análisis estructural de resonancias magnéticas lumbares, dado que la segmentación de estructuras anatómicas, la obtención de mediciones simples y la generación de salidas revisables aún pueden requerir tareas manuales, variables y poco integradas dentro de un mismo flujo funcional.

La dificultad no se ubica en reemplazar la interpretación profesional, sino en asistir tareas operativas específicas. Un sistema de apoyo puede proponer segmentaciones, calcular mediciones geométricas y presentar resultados visuales, pero debe permitir revisión, corrección o descarte por parte del usuario. Desde una perspectiva de ingeniería, el desafío consiste en integrar carga de imágenes, preprocesamiento, segmentación automática, mediciones, visualización superpuesta y salida editable en un flujo coherente. La solución debe ser trazable, reproducible y acotada al contexto académico, sin presentarse como producto clínico certificado ni como sistema de diagnóstico autónomo.
